%% Deterministic equivalent for Var(R) — the second-order correction to
%% the accentuation error.  Import into main document via %% Deterministic equivalent for Var(R) — the second-order correction to
%% the accentuation error.  Import into main document via %% Deterministic equivalent for Var(R) — the second-order correction to
%% the accentuation error.  Import into main document via %% Deterministic equivalent for Var(R) — the second-order correction to
%% the accentuation error.  Import into main document via \input{accentuation_var_correction}.

\subsection{Second-order correction to the accentuation error: $\operatorname{Var}(R)$}

\paragraph{Why the leading-order formula is insufficient.}

The accentuation error derived above relies on the approximation
\begin{equation}
  \mathbb{E}_{X,\eta}\!\left[\left(1-R\right)^{2}\right]
  \;\approx\;
  \left(1-R_{\det}\right)^{2},
  \label{eq:acc_leading}
\end{equation}
where $R = \hat{\beta}^{\top}\beta^{*}/\hat{\beta}^{\top}\hat{\beta}$ and
$R_{\det}$ is its deterministic equivalent.  By the second-moment
decomposition,
\begin{equation}
  \mathbb{E}\!\left[(1-R)^{2}\right]
  \;=\;
  \underbrace{(1-\mathbb{E}[R])^{2}}_{\text{bias}^{2}}
  \;+\;
  \underbrace{\operatorname{Var}(R)}_{\text{missing term}}.
  \label{eq:var_decomp}
\end{equation}
The leading-order theory captures $(1-R_{\det})^{2}\approx(1-\mathbb{E}[R])^{2}$
but ignores $\operatorname{Var}(R)$.  When $R\approx 1$
(well-regularized regime), the bias$^{2}$ term is tiny while
$\operatorname{Var}(R)=O(1/n)$ remains finite, so $\operatorname{Var}(R)$
can dominate.

\paragraph{Delta-method setup.}

Define the numerator and denominator of $R$:
\[
  N \;=\; \hat{\beta}^{\top}\beta^{*},
  \qquad
  D \;=\; \hat{\beta}^{\top}\hat{\beta}.
\]
Both $N$ and $D$ concentrate in the large-$n$, fixed-$\gamma=d/n$ limit:
$N\to N_{\det}$, $D\to D_{\det}$, with fluctuations of order $1/\sqrt{n}$.
By the delta method,
\begin{equation}
  R - R_{\det}
  \;\approx\;
  \frac{1}{D_{\det}}\left[(N-N_{\det}) - R_{\det}(D-D_{\det})\right]
  \;=\;
  \frac{\delta Q}{D_{\det}},
  \label{eq:deltaR}
\end{equation}
where $\delta Q = (N - R_{\det}D) - \mathbb{E}[N-R_{\det}D]$.
Because $R_{\det}=N_{\det}/D_{\det}$, the mean of $Q=N-R_{\det}D$ is
zero by construction, so $\delta Q = N - R_{\det}D$ exactly.
Therefore,
\begin{equation}
  \operatorname{Var}(R)
  \;\asymp\;
  \frac{\operatorname{Var}(Q)}{D_{\det}^{2}},
  \quad
  Q \;=\; \hat{\beta}^{\top}(\beta^{*} - R_{\det}\hat{\beta}).
  \label{eq:varR_via_Q}
\end{equation}

\paragraph{Decomposing $\hat{\beta}$ into signal and noise.}

Write $\hat{\beta} = \mathbf{g} + \mathbf{h}$ where
\begin{align*}
  \mathbf{g} &\;=\; (\hat{\Sigma}+\lambda I)^{-1}\hat{\Sigma}\beta^{*}
  \;\asymp\;(\Sigma+\kappa I)^{-1}\Sigma\beta^{*},\\
  \mathbf{h} &\;=\; \sigma(\hat{\Sigma}+\lambda I)^{-1}\!\left(\tfrac{1}{n}X^{\top}\eta\right)
  \;\sim\;
  \mathcal{N}\!\left(0,\,\tfrac{\sigma^{2}}{n}H\right),
  \quad
  H = (\Sigma+\kappa I)^{-1}\Sigma(\Sigma+\kappa I)^{-1}.
\end{align*}
Then
\[
  Q \;=\; (\mathbf{g}+\mathbf{h})^{\top}\bigl(\beta^{*}-R_{\det}(\mathbf{g}+\mathbf{h})\bigr).
\]
Expanding and using $N_{\det}-R_{\det}D_{\det}=0$,
\begin{equation}
  Q
  \;=\;
  \underbrace{\mathbf{h}^{\top}\bigl(\beta^{*}-2R_{\det}\mathbf{g}\bigr)}_{\delta Q^{(\mathrm{lin})}}
  \;-\;
  R_{\det}\!\underbrace{\bigl(\mathbf{h}^{\top}\mathbf{h} - \mathbb{E}[\mathbf{h}^{\top}\mathbf{h}]\bigr)}_{\delta Q^{(\mathrm{quad})}}.
  \label{eq:Q_expand}
\end{equation}
The two terms have zero mean and are uncorrelated (parity argument:
$\delta Q^{(\mathrm{lin})}$ is odd in $\mathbf{h}$, $\delta Q^{(\mathrm{quad})}$
is even), so $\operatorname{Var}(Q)=\operatorname{Var}(\delta Q^{(\mathrm{lin})})+\operatorname{Var}(\delta Q^{(\mathrm{quad})})$.

\paragraph{Linear (noise) contribution.}

Since $\mathbf{h}\sim\mathcal{N}(0,\frac{\sigma^{2}}{n}H)$,
\[
  \operatorname{Var}(\delta Q^{(\mathrm{lin})})
  \;=\;
  \frac{\sigma^{2}}{n}
  \bigl(\beta^{*}-2R_{\det}\mathbf{g}_{\det}\bigr)^{\top}H\,
  \bigl(\beta^{*}-2R_{\det}\mathbf{g}_{\det}\bigr),
\]
with $\mathbf{g}_{\det}=(\Sigma+\kappa I)^{-1}\Sigma\beta^{*}$.
In the eigenbasis of $\Sigma$ (eigenvalues $\lambda_{k}$,
projections $b_{k}=\mathbf{u}_{k}^{\top}\beta^{*}$):
\[
  H_{kk} = \frac{\lambda_{k}}{(\lambda_{k}+\kappa)^{2}},
  \quad
  \bigl(g_{\det}\bigr)_{k} = \frac{\lambda_{k}}{\lambda_{k}+\kappa}b_{k},
  \quad
  w_{k} := \left(\beta^{*}-2R_{\det}\mathbf{g}_{\det}\right)_{k}
  = b_{k}\!\left(1 - \frac{2R_{\det}\lambda_{k}}{\lambda_{k}+\kappa}\right).
\]
Hence,
\begin{equation}
  \operatorname{Var}(\delta Q^{(\mathrm{lin})})
  \;=\;
  \frac{\sigma^{2}}{n}\sum_{k}\frac{\lambda_{k}b_{k}^{2}}{(\lambda_{k}+\kappa)^{2}}
  \!\left(1-\frac{2R_{\det}\lambda_{k}}{\lambda_{k}+\kappa}\right)^{\!2}.
  \label{eq:varQ_lin}
\end{equation}

Expanding the square and introducing shorthand sums:
\begin{align}
  C_{\mathrm{sig}} &:= \sum_{k}\frac{\lambda_{k}}{(\lambda_{k}+\kappa)^{2}}b_{k}^{2},
  \label{eq:Csig}\\
  M_{1} &:= \sum_{k}\frac{\lambda_{k}^{2}}{(\lambda_{k}+\kappa)^{3}}b_{k}^{2},
  \label{eq:M1}\\
  M_{2} &:= \sum_{k}\frac{\lambda_{k}^{3}}{(\lambda_{k}+\kappa)^{4}}b_{k}^{2},
  \label{eq:M2}
\end{align}
equation~\eqref{eq:varQ_lin} becomes
\begin{equation}
  \operatorname{Var}(\delta Q^{(\mathrm{lin})})
  \;=\;
  \frac{\sigma^{2}}{n}\Bigl[C_{\mathrm{sig}} - 4R_{\det}M_{1} + 4R_{\det}^{2}M_{2}\Bigr].
  \label{eq:varQ_lin_compact}
\end{equation}

\paragraph{Quadratic correction.}

$\delta Q^{(\mathrm{quad})} = \mathbf{h}^{\top}\mathbf{h}-\mathbb{E}[\mathbf{h}^{\top}\mathbf{h}]$ is
a centered chi-squared form.  Its variance is
\[
  \operatorname{Var}(\delta Q^{(\mathrm{quad})})
  \;=\;
  2\left(\frac{\sigma^{2}}{n}\right)^{\!2}\operatorname{Tr}[H^{2}]
  \;=\;
  \frac{2\sigma^{4}}{n^{2}}\sum_{k}\frac{\lambda_{k}^{2}}{(\lambda_{k}+\kappa)^{4}}.
  \label{eq:varQ_quad}
\]
This term is $O(\sigma^{4}\gamma/n)$ and dominates when $\sigma^{2}\gg 1/\sqrt{\gamma}$,
i.e.\ at high noise with fixed aspect ratio.

\paragraph{Deterministic equivalent for $\operatorname{Var}(R)$.}

Combining equations~\eqref{eq:varR_via_Q},~\eqref{eq:varQ_lin_compact},
and~\eqref{eq:varQ_quad}:

\begin{equation}
  \boxed{
  \operatorname{Var}(R)
  \;\asymp\;
  \frac{1}{D_{\det}^{2}}
  \left[
    \frac{\sigma^{2}}{n}\Bigl(C_{\mathrm{sig}} - 4R_{\det}M_{1} + 4R_{\det}^{2}M_{2}\Bigr)
    \;+\;
    \frac{2R_{\det}^{2}\sigma^{4}}{n^{2}}\sum_{k}\frac{\lambda_{k}^{2}}{(\lambda_{k}+\kappa)^{4}}
  \right]
  }
  \label{eq:varR_det}
\end{equation}

\paragraph{Full accentuation error including $\operatorname{Var}(R)$.}

The improved deterministic equivalent for the accentuation error is:
\begin{equation}
  \mathbb{E}\!\left[E_{\mathrm{acc}}\right]
  \;\asymp\;
  \beta^{*\top}\Sigma\beta^{*}
  \cdot
  \left[
    (1-R_{\det})^{2}
    \;+\;
    \operatorname{Var}(R)
  \right],
  \label{eq:Eacc_full}
\end{equation}
where $\operatorname{Var}(R)$ is given by~\eqref{eq:varR_det}.

\paragraph{Regime analysis.}
\begin{itemize}
\item \textbf{Low noise} ($\sigma\to 0$): the $\eta$-driven terms
  in~\eqref{eq:varR_det} vanish as $O(\sigma^{2})$, but $\operatorname{Var}(R)$
  remains finite due to \emph{X-fluctuations}
  $\operatorname{Var}_{X}(\mathbb{E}_{\eta}[R|X])=O(1/n)$.
  At low $\sigma$, X-fluctuations dominate
  (empirically $\approx 90\%$ of total $\operatorname{Var}(R)$) and are
  \emph{not} captured by equation~\eqref{eq:varR_det}.
  They require a CLT for quadratic functionals of Wishart matrices
  (a higher-order RMT result).
\item \textbf{Moderate-to-high noise}: the $\eta$-driven terms grow as
  $O(\sigma^{2}/n)$ and quickly dominate ($\approx 97\%$ at $\sigma=1$).
  The quadratic term $O(\sigma^{4}/n)$ can exceed the linear term when
  $\sigma^{2}\sqrt{\gamma}\gg 1$.
\item \textbf{Det.\ equiv.\ accuracy}: even within the $\eta$-driven
  contribution, replacing $H(X)$ by $H_{\det}$ introduces systematic
  error.  When $\kappa=O(\lambda)$ is small relative to the eigenvalue
  spread, the det.\ equiv.\ under-approximates $\operatorname{Tr}[H^{2}]$
  and $\mathbf{w}^{\top}H\mathbf{w}$ (ratio $\sim 0.4$--$0.7$
  in the van Hateren example at $\lambda^{*}\approx 0.001$).
  Using the actual sample eigenvalues of $\hat{\Sigma}$ in place of
  the population ones restores the accuracy (ratio $\to\approx 0.95$).
\end{itemize}

\paragraph{Empirical decomposition (van Hateren $16\times 16$, $d=256$, $n=512$).}

The law of total variance $\operatorname{Var}(R)=\mathbb{E}_{X}[\operatorname{Var}_{\eta}(R|X)]+\operatorname{Var}_{X}(\mathbb{E}_{\eta}[R|X])$ was measured with $T=300$ datasets $\times$ $T'=80$ noise replicas each:

\begin{center}
\begin{tabular}{ccccc}
\hline
$\sigma$ & $\lambda^{*}$ & $\eta$-fraction & X-fraction & formula / $\eta$-MC \\
\hline
0.10 & 0.001 & 10\% & 90\% & 0.69 \\
1.00 & 0.002 & 98\% & 2\%  & 0.44 \\
\hline
\end{tabular}
\end{center}

The formula captures the correct order of magnitude and correctly predicts
that $\operatorname{Var}(R)\sim O(1/n)$, but a complete first-principles
derivation of $\operatorname{Var}(R)$ requires accounting for both
X-fluctuations and the finite-$\kappa$ correction to the det.\ equiv.\
of $H$.
.

\subsection{Second-order correction to the accentuation error: $\operatorname{Var}(R)$}

\paragraph{Why the leading-order formula is insufficient.}

The accentuation error derived above relies on the approximation
\begin{equation}
  \mathbb{E}_{X,\eta}\!\left[\left(1-R\right)^{2}\right]
  \;\approx\;
  \left(1-R_{\det}\right)^{2},
  \label{eq:acc_leading}
\end{equation}
where $R = \hat{\beta}^{\top}\beta^{*}/\hat{\beta}^{\top}\hat{\beta}$ and
$R_{\det}$ is its deterministic equivalent.  By the second-moment
decomposition,
\begin{equation}
  \mathbb{E}\!\left[(1-R)^{2}\right]
  \;=\;
  \underbrace{(1-\mathbb{E}[R])^{2}}_{\text{bias}^{2}}
  \;+\;
  \underbrace{\operatorname{Var}(R)}_{\text{missing term}}.
  \label{eq:var_decomp}
\end{equation}
The leading-order theory captures $(1-R_{\det})^{2}\approx(1-\mathbb{E}[R])^{2}$
but ignores $\operatorname{Var}(R)$.  When $R\approx 1$
(well-regularized regime), the bias$^{2}$ term is tiny while
$\operatorname{Var}(R)=O(1/n)$ remains finite, so $\operatorname{Var}(R)$
can dominate.

\paragraph{Delta-method setup.}

Define the numerator and denominator of $R$:
\[
  N \;=\; \hat{\beta}^{\top}\beta^{*},
  \qquad
  D \;=\; \hat{\beta}^{\top}\hat{\beta}.
\]
Both $N$ and $D$ concentrate in the large-$n$, fixed-$\gamma=d/n$ limit:
$N\to N_{\det}$, $D\to D_{\det}$, with fluctuations of order $1/\sqrt{n}$.
By the delta method,
\begin{equation}
  R - R_{\det}
  \;\approx\;
  \frac{1}{D_{\det}}\left[(N-N_{\det}) - R_{\det}(D-D_{\det})\right]
  \;=\;
  \frac{\delta Q}{D_{\det}},
  \label{eq:deltaR}
\end{equation}
where $\delta Q = (N - R_{\det}D) - \mathbb{E}[N-R_{\det}D]$.
Because $R_{\det}=N_{\det}/D_{\det}$, the mean of $Q=N-R_{\det}D$ is
zero by construction, so $\delta Q = N - R_{\det}D$ exactly.
Therefore,
\begin{equation}
  \operatorname{Var}(R)
  \;\asymp\;
  \frac{\operatorname{Var}(Q)}{D_{\det}^{2}},
  \quad
  Q \;=\; \hat{\beta}^{\top}(\beta^{*} - R_{\det}\hat{\beta}).
  \label{eq:varR_via_Q}
\end{equation}

\paragraph{Decomposing $\hat{\beta}$ into signal and noise.}

Write $\hat{\beta} = \mathbf{g} + \mathbf{h}$ where
\begin{align*}
  \mathbf{g} &\;=\; (\hat{\Sigma}+\lambda I)^{-1}\hat{\Sigma}\beta^{*}
  \;\asymp\;(\Sigma+\kappa I)^{-1}\Sigma\beta^{*},\\
  \mathbf{h} &\;=\; \sigma(\hat{\Sigma}+\lambda I)^{-1}\!\left(\tfrac{1}{n}X^{\top}\eta\right)
  \;\sim\;
  \mathcal{N}\!\left(0,\,\tfrac{\sigma^{2}}{n}H\right),
  \quad
  H = (\Sigma+\kappa I)^{-1}\Sigma(\Sigma+\kappa I)^{-1}.
\end{align*}
Then
\[
  Q \;=\; (\mathbf{g}+\mathbf{h})^{\top}\bigl(\beta^{*}-R_{\det}(\mathbf{g}+\mathbf{h})\bigr).
\]
Expanding and using $N_{\det}-R_{\det}D_{\det}=0$,
\begin{equation}
  Q
  \;=\;
  \underbrace{\mathbf{h}^{\top}\bigl(\beta^{*}-2R_{\det}\mathbf{g}\bigr)}_{\delta Q^{(\mathrm{lin})}}
  \;-\;
  R_{\det}\!\underbrace{\bigl(\mathbf{h}^{\top}\mathbf{h} - \mathbb{E}[\mathbf{h}^{\top}\mathbf{h}]\bigr)}_{\delta Q^{(\mathrm{quad})}}.
  \label{eq:Q_expand}
\end{equation}
The two terms have zero mean and are uncorrelated (parity argument:
$\delta Q^{(\mathrm{lin})}$ is odd in $\mathbf{h}$, $\delta Q^{(\mathrm{quad})}$
is even), so $\operatorname{Var}(Q)=\operatorname{Var}(\delta Q^{(\mathrm{lin})})+\operatorname{Var}(\delta Q^{(\mathrm{quad})})$.

\paragraph{Linear (noise) contribution.}

Since $\mathbf{h}\sim\mathcal{N}(0,\frac{\sigma^{2}}{n}H)$,
\[
  \operatorname{Var}(\delta Q^{(\mathrm{lin})})
  \;=\;
  \frac{\sigma^{2}}{n}
  \bigl(\beta^{*}-2R_{\det}\mathbf{g}_{\det}\bigr)^{\top}H\,
  \bigl(\beta^{*}-2R_{\det}\mathbf{g}_{\det}\bigr),
\]
with $\mathbf{g}_{\det}=(\Sigma+\kappa I)^{-1}\Sigma\beta^{*}$.
In the eigenbasis of $\Sigma$ (eigenvalues $\lambda_{k}$,
projections $b_{k}=\mathbf{u}_{k}^{\top}\beta^{*}$):
\[
  H_{kk} = \frac{\lambda_{k}}{(\lambda_{k}+\kappa)^{2}},
  \quad
  \bigl(g_{\det}\bigr)_{k} = \frac{\lambda_{k}}{\lambda_{k}+\kappa}b_{k},
  \quad
  w_{k} := \left(\beta^{*}-2R_{\det}\mathbf{g}_{\det}\right)_{k}
  = b_{k}\!\left(1 - \frac{2R_{\det}\lambda_{k}}{\lambda_{k}+\kappa}\right).
\]
Hence,
\begin{equation}
  \operatorname{Var}(\delta Q^{(\mathrm{lin})})
  \;=\;
  \frac{\sigma^{2}}{n}\sum_{k}\frac{\lambda_{k}b_{k}^{2}}{(\lambda_{k}+\kappa)^{2}}
  \!\left(1-\frac{2R_{\det}\lambda_{k}}{\lambda_{k}+\kappa}\right)^{\!2}.
  \label{eq:varQ_lin}
\end{equation}

Expanding the square and introducing shorthand sums:
\begin{align}
  C_{\mathrm{sig}} &:= \sum_{k}\frac{\lambda_{k}}{(\lambda_{k}+\kappa)^{2}}b_{k}^{2},
  \label{eq:Csig}\\
  M_{1} &:= \sum_{k}\frac{\lambda_{k}^{2}}{(\lambda_{k}+\kappa)^{3}}b_{k}^{2},
  \label{eq:M1}\\
  M_{2} &:= \sum_{k}\frac{\lambda_{k}^{3}}{(\lambda_{k}+\kappa)^{4}}b_{k}^{2},
  \label{eq:M2}
\end{align}
equation~\eqref{eq:varQ_lin} becomes
\begin{equation}
  \operatorname{Var}(\delta Q^{(\mathrm{lin})})
  \;=\;
  \frac{\sigma^{2}}{n}\Bigl[C_{\mathrm{sig}} - 4R_{\det}M_{1} + 4R_{\det}^{2}M_{2}\Bigr].
  \label{eq:varQ_lin_compact}
\end{equation}

\paragraph{Quadratic correction.}

$\delta Q^{(\mathrm{quad})} = \mathbf{h}^{\top}\mathbf{h}-\mathbb{E}[\mathbf{h}^{\top}\mathbf{h}]$ is
a centered chi-squared form.  Its variance is
\[
  \operatorname{Var}(\delta Q^{(\mathrm{quad})})
  \;=\;
  2\left(\frac{\sigma^{2}}{n}\right)^{\!2}\operatorname{Tr}[H^{2}]
  \;=\;
  \frac{2\sigma^{4}}{n^{2}}\sum_{k}\frac{\lambda_{k}^{2}}{(\lambda_{k}+\kappa)^{4}}.
  \label{eq:varQ_quad}
\]
This term is $O(\sigma^{4}\gamma/n)$ and dominates when $\sigma^{2}\gg 1/\sqrt{\gamma}$,
i.e.\ at high noise with fixed aspect ratio.

\paragraph{Deterministic equivalent for $\operatorname{Var}(R)$.}

Combining equations~\eqref{eq:varR_via_Q},~\eqref{eq:varQ_lin_compact},
and~\eqref{eq:varQ_quad}:

\begin{equation}
  \boxed{
  \operatorname{Var}(R)
  \;\asymp\;
  \frac{1}{D_{\det}^{2}}
  \left[
    \frac{\sigma^{2}}{n}\Bigl(C_{\mathrm{sig}} - 4R_{\det}M_{1} + 4R_{\det}^{2}M_{2}\Bigr)
    \;+\;
    \frac{2R_{\det}^{2}\sigma^{4}}{n^{2}}\sum_{k}\frac{\lambda_{k}^{2}}{(\lambda_{k}+\kappa)^{4}}
  \right]
  }
  \label{eq:varR_det}
\end{equation}

\paragraph{Full accentuation error including $\operatorname{Var}(R)$.}

The improved deterministic equivalent for the accentuation error is:
\begin{equation}
  \mathbb{E}\!\left[E_{\mathrm{acc}}\right]
  \;\asymp\;
  \beta^{*\top}\Sigma\beta^{*}
  \cdot
  \left[
    (1-R_{\det})^{2}
    \;+\;
    \operatorname{Var}(R)
  \right],
  \label{eq:Eacc_full}
\end{equation}
where $\operatorname{Var}(R)$ is given by~\eqref{eq:varR_det}.

\paragraph{Regime analysis.}
\begin{itemize}
\item \textbf{Low noise} ($\sigma\to 0$): the $\eta$-driven terms
  in~\eqref{eq:varR_det} vanish as $O(\sigma^{2})$, but $\operatorname{Var}(R)$
  remains finite due to \emph{X-fluctuations}
  $\operatorname{Var}_{X}(\mathbb{E}_{\eta}[R|X])=O(1/n)$.
  At low $\sigma$, X-fluctuations dominate
  (empirically $\approx 90\%$ of total $\operatorname{Var}(R)$) and are
  \emph{not} captured by equation~\eqref{eq:varR_det}.
  They require a CLT for quadratic functionals of Wishart matrices
  (a higher-order RMT result).
\item \textbf{Moderate-to-high noise}: the $\eta$-driven terms grow as
  $O(\sigma^{2}/n)$ and quickly dominate ($\approx 97\%$ at $\sigma=1$).
  The quadratic term $O(\sigma^{4}/n)$ can exceed the linear term when
  $\sigma^{2}\sqrt{\gamma}\gg 1$.
\item \textbf{Det.\ equiv.\ accuracy}: even within the $\eta$-driven
  contribution, replacing $H(X)$ by $H_{\det}$ introduces systematic
  error.  When $\kappa=O(\lambda)$ is small relative to the eigenvalue
  spread, the det.\ equiv.\ under-approximates $\operatorname{Tr}[H^{2}]$
  and $\mathbf{w}^{\top}H\mathbf{w}$ (ratio $\sim 0.4$--$0.7$
  in the van Hateren example at $\lambda^{*}\approx 0.001$).
  Using the actual sample eigenvalues of $\hat{\Sigma}$ in place of
  the population ones restores the accuracy (ratio $\to\approx 0.95$).
\end{itemize}

\paragraph{Empirical decomposition (van Hateren $16\times 16$, $d=256$, $n=512$).}

The law of total variance $\operatorname{Var}(R)=\mathbb{E}_{X}[\operatorname{Var}_{\eta}(R|X)]+\operatorname{Var}_{X}(\mathbb{E}_{\eta}[R|X])$ was measured with $T=300$ datasets $\times$ $T'=80$ noise replicas each:

\begin{center}
\begin{tabular}{ccccc}
\hline
$\sigma$ & $\lambda^{*}$ & $\eta$-fraction & X-fraction & formula / $\eta$-MC \\
\hline
0.10 & 0.001 & 10\% & 90\% & 0.69 \\
1.00 & 0.002 & 98\% & 2\%  & 0.44 \\
\hline
\end{tabular}
\end{center}

The formula captures the correct order of magnitude and correctly predicts
that $\operatorname{Var}(R)\sim O(1/n)$, but a complete first-principles
derivation of $\operatorname{Var}(R)$ requires accounting for both
X-fluctuations and the finite-$\kappa$ correction to the det.\ equiv.\
of $H$.
.

\subsection{Second-order correction to the accentuation error: $\operatorname{Var}(R)$}

\paragraph{Why the leading-order formula is insufficient.}

The accentuation error derived above relies on the approximation
\begin{equation}
  \mathbb{E}_{X,\eta}\!\left[\left(1-R\right)^{2}\right]
  \;\approx\;
  \left(1-R_{\det}\right)^{2},
  \label{eq:acc_leading}
\end{equation}
where $R = \hat{\beta}^{\top}\beta^{*}/\hat{\beta}^{\top}\hat{\beta}$ and
$R_{\det}$ is its deterministic equivalent.  By the second-moment
decomposition,
\begin{equation}
  \mathbb{E}\!\left[(1-R)^{2}\right]
  \;=\;
  \underbrace{(1-\mathbb{E}[R])^{2}}_{\text{bias}^{2}}
  \;+\;
  \underbrace{\operatorname{Var}(R)}_{\text{missing term}}.
  \label{eq:var_decomp}
\end{equation}
The leading-order theory captures $(1-R_{\det})^{2}\approx(1-\mathbb{E}[R])^{2}$
but ignores $\operatorname{Var}(R)$.  When $R\approx 1$
(well-regularized regime), the bias$^{2}$ term is tiny while
$\operatorname{Var}(R)=O(1/n)$ remains finite, so $\operatorname{Var}(R)$
can dominate.

\paragraph{Delta-method setup.}

Define the numerator and denominator of $R$:
\[
  N \;=\; \hat{\beta}^{\top}\beta^{*},
  \qquad
  D \;=\; \hat{\beta}^{\top}\hat{\beta}.
\]
Both $N$ and $D$ concentrate in the large-$n$, fixed-$\gamma=d/n$ limit:
$N\to N_{\det}$, $D\to D_{\det}$, with fluctuations of order $1/\sqrt{n}$.
By the delta method,
\begin{equation}
  R - R_{\det}
  \;\approx\;
  \frac{1}{D_{\det}}\left[(N-N_{\det}) - R_{\det}(D-D_{\det})\right]
  \;=\;
  \frac{\delta Q}{D_{\det}},
  \label{eq:deltaR}
\end{equation}
where $\delta Q = (N - R_{\det}D) - \mathbb{E}[N-R_{\det}D]$.
Because $R_{\det}=N_{\det}/D_{\det}$, the mean of $Q=N-R_{\det}D$ is
zero by construction, so $\delta Q = N - R_{\det}D$ exactly.
Therefore,
\begin{equation}
  \operatorname{Var}(R)
  \;\asymp\;
  \frac{\operatorname{Var}(Q)}{D_{\det}^{2}},
  \quad
  Q \;=\; \hat{\beta}^{\top}(\beta^{*} - R_{\det}\hat{\beta}).
  \label{eq:varR_via_Q}
\end{equation}

\paragraph{Decomposing $\hat{\beta}$ into signal and noise.}

Write $\hat{\beta} = \mathbf{g} + \mathbf{h}$ where
\begin{align*}
  \mathbf{g} &\;=\; (\hat{\Sigma}+\lambda I)^{-1}\hat{\Sigma}\beta^{*}
  \;\asymp\;(\Sigma+\kappa I)^{-1}\Sigma\beta^{*},\\
  \mathbf{h} &\;=\; \sigma(\hat{\Sigma}+\lambda I)^{-1}\!\left(\tfrac{1}{n}X^{\top}\eta\right)
  \;\sim\;
  \mathcal{N}\!\left(0,\,\tfrac{\sigma^{2}}{n}H\right),
  \quad
  H = (\Sigma+\kappa I)^{-1}\Sigma(\Sigma+\kappa I)^{-1}.
\end{align*}
Then
\[
  Q \;=\; (\mathbf{g}+\mathbf{h})^{\top}\bigl(\beta^{*}-R_{\det}(\mathbf{g}+\mathbf{h})\bigr).
\]
Expanding and using $N_{\det}-R_{\det}D_{\det}=0$,
\begin{equation}
  Q
  \;=\;
  \underbrace{\mathbf{h}^{\top}\bigl(\beta^{*}-2R_{\det}\mathbf{g}\bigr)}_{\delta Q^{(\mathrm{lin})}}
  \;-\;
  R_{\det}\!\underbrace{\bigl(\mathbf{h}^{\top}\mathbf{h} - \mathbb{E}[\mathbf{h}^{\top}\mathbf{h}]\bigr)}_{\delta Q^{(\mathrm{quad})}}.
  \label{eq:Q_expand}
\end{equation}
The two terms have zero mean and are uncorrelated (parity argument:
$\delta Q^{(\mathrm{lin})}$ is odd in $\mathbf{h}$, $\delta Q^{(\mathrm{quad})}$
is even), so $\operatorname{Var}(Q)=\operatorname{Var}(\delta Q^{(\mathrm{lin})})+\operatorname{Var}(\delta Q^{(\mathrm{quad})})$.

\paragraph{Linear (noise) contribution.}

Since $\mathbf{h}\sim\mathcal{N}(0,\frac{\sigma^{2}}{n}H)$,
\[
  \operatorname{Var}(\delta Q^{(\mathrm{lin})})
  \;=\;
  \frac{\sigma^{2}}{n}
  \bigl(\beta^{*}-2R_{\det}\mathbf{g}_{\det}\bigr)^{\top}H\,
  \bigl(\beta^{*}-2R_{\det}\mathbf{g}_{\det}\bigr),
\]
with $\mathbf{g}_{\det}=(\Sigma+\kappa I)^{-1}\Sigma\beta^{*}$.
In the eigenbasis of $\Sigma$ (eigenvalues $\lambda_{k}$,
projections $b_{k}=\mathbf{u}_{k}^{\top}\beta^{*}$):
\[
  H_{kk} = \frac{\lambda_{k}}{(\lambda_{k}+\kappa)^{2}},
  \quad
  \bigl(g_{\det}\bigr)_{k} = \frac{\lambda_{k}}{\lambda_{k}+\kappa}b_{k},
  \quad
  w_{k} := \left(\beta^{*}-2R_{\det}\mathbf{g}_{\det}\right)_{k}
  = b_{k}\!\left(1 - \frac{2R_{\det}\lambda_{k}}{\lambda_{k}+\kappa}\right).
\]
Hence,
\begin{equation}
  \operatorname{Var}(\delta Q^{(\mathrm{lin})})
  \;=\;
  \frac{\sigma^{2}}{n}\sum_{k}\frac{\lambda_{k}b_{k}^{2}}{(\lambda_{k}+\kappa)^{2}}
  \!\left(1-\frac{2R_{\det}\lambda_{k}}{\lambda_{k}+\kappa}\right)^{\!2}.
  \label{eq:varQ_lin}
\end{equation}

Expanding the square and introducing shorthand sums:
\begin{align}
  C_{\mathrm{sig}} &:= \sum_{k}\frac{\lambda_{k}}{(\lambda_{k}+\kappa)^{2}}b_{k}^{2},
  \label{eq:Csig}\\
  M_{1} &:= \sum_{k}\frac{\lambda_{k}^{2}}{(\lambda_{k}+\kappa)^{3}}b_{k}^{2},
  \label{eq:M1}\\
  M_{2} &:= \sum_{k}\frac{\lambda_{k}^{3}}{(\lambda_{k}+\kappa)^{4}}b_{k}^{2},
  \label{eq:M2}
\end{align}
equation~\eqref{eq:varQ_lin} becomes
\begin{equation}
  \operatorname{Var}(\delta Q^{(\mathrm{lin})})
  \;=\;
  \frac{\sigma^{2}}{n}\Bigl[C_{\mathrm{sig}} - 4R_{\det}M_{1} + 4R_{\det}^{2}M_{2}\Bigr].
  \label{eq:varQ_lin_compact}
\end{equation}

\paragraph{Quadratic correction.}

$\delta Q^{(\mathrm{quad})} = \mathbf{h}^{\top}\mathbf{h}-\mathbb{E}[\mathbf{h}^{\top}\mathbf{h}]$ is
a centered chi-squared form.  Its variance is
\[
  \operatorname{Var}(\delta Q^{(\mathrm{quad})})
  \;=\;
  2\left(\frac{\sigma^{2}}{n}\right)^{\!2}\operatorname{Tr}[H^{2}]
  \;=\;
  \frac{2\sigma^{4}}{n^{2}}\sum_{k}\frac{\lambda_{k}^{2}}{(\lambda_{k}+\kappa)^{4}}.
  \label{eq:varQ_quad}
\]
This term is $O(\sigma^{4}\gamma/n)$ and dominates when $\sigma^{2}\gg 1/\sqrt{\gamma}$,
i.e.\ at high noise with fixed aspect ratio.

\paragraph{Deterministic equivalent for $\operatorname{Var}(R)$.}

Combining equations~\eqref{eq:varR_via_Q},~\eqref{eq:varQ_lin_compact},
and~\eqref{eq:varQ_quad}:

\begin{equation}
  \boxed{
  \operatorname{Var}(R)
  \;\asymp\;
  \frac{1}{D_{\det}^{2}}
  \left[
    \frac{\sigma^{2}}{n}\Bigl(C_{\mathrm{sig}} - 4R_{\det}M_{1} + 4R_{\det}^{2}M_{2}\Bigr)
    \;+\;
    \frac{2R_{\det}^{2}\sigma^{4}}{n^{2}}\sum_{k}\frac{\lambda_{k}^{2}}{(\lambda_{k}+\kappa)^{4}}
  \right]
  }
  \label{eq:varR_det}
\end{equation}

\paragraph{Full accentuation error including $\operatorname{Var}(R)$.}

The improved deterministic equivalent for the accentuation error is:
\begin{equation}
  \mathbb{E}\!\left[E_{\mathrm{acc}}\right]
  \;\asymp\;
  \beta^{*\top}\Sigma\beta^{*}
  \cdot
  \left[
    (1-R_{\det})^{2}
    \;+\;
    \operatorname{Var}(R)
  \right],
  \label{eq:Eacc_full}
\end{equation}
where $\operatorname{Var}(R)$ is given by~\eqref{eq:varR_det}.

\paragraph{Regime analysis.}
\begin{itemize}
\item \textbf{Low noise} ($\sigma\to 0$): the $\eta$-driven terms
  in~\eqref{eq:varR_det} vanish as $O(\sigma^{2})$, but $\operatorname{Var}(R)$
  remains finite due to \emph{X-fluctuations}
  $\operatorname{Var}_{X}(\mathbb{E}_{\eta}[R|X])=O(1/n)$.
  At low $\sigma$, X-fluctuations dominate
  (empirically $\approx 90\%$ of total $\operatorname{Var}(R)$) and are
  \emph{not} captured by equation~\eqref{eq:varR_det}.
  They require a CLT for quadratic functionals of Wishart matrices
  (a higher-order RMT result).
\item \textbf{Moderate-to-high noise}: the $\eta$-driven terms grow as
  $O(\sigma^{2}/n)$ and quickly dominate ($\approx 97\%$ at $\sigma=1$).
  The quadratic term $O(\sigma^{4}/n)$ can exceed the linear term when
  $\sigma^{2}\sqrt{\gamma}\gg 1$.
\item \textbf{Det.\ equiv.\ accuracy}: even within the $\eta$-driven
  contribution, replacing $H(X)$ by $H_{\det}$ introduces systematic
  error.  When $\kappa=O(\lambda)$ is small relative to the eigenvalue
  spread, the det.\ equiv.\ under-approximates $\operatorname{Tr}[H^{2}]$
  and $\mathbf{w}^{\top}H\mathbf{w}$ (ratio $\sim 0.4$--$0.7$
  in the van Hateren example at $\lambda^{*}\approx 0.001$).
  Using the actual sample eigenvalues of $\hat{\Sigma}$ in place of
  the population ones restores the accuracy (ratio $\to\approx 0.95$).
\end{itemize}

\paragraph{Empirical decomposition (van Hateren $16\times 16$, $d=256$, $n=512$).}

The law of total variance $\operatorname{Var}(R)=\mathbb{E}_{X}[\operatorname{Var}_{\eta}(R|X)]+\operatorname{Var}_{X}(\mathbb{E}_{\eta}[R|X])$ was measured with $T=300$ datasets $\times$ $T'=80$ noise replicas each:

\begin{center}
\begin{tabular}{ccccc}
\hline
$\sigma$ & $\lambda^{*}$ & $\eta$-fraction & X-fraction & formula / $\eta$-MC \\
\hline
0.10 & 0.001 & 10\% & 90\% & 0.69 \\
1.00 & 0.002 & 98\% & 2\%  & 0.44 \\
\hline
\end{tabular}
\end{center}

The formula captures the correct order of magnitude and correctly predicts
that $\operatorname{Var}(R)\sim O(1/n)$, but a complete first-principles
derivation of $\operatorname{Var}(R)$ requires accounting for both
X-fluctuations and the finite-$\kappa$ correction to the det.\ equiv.\
of $H$.
.

\subsection{Second-order correction to the accentuation error: $\operatorname{Var}(R)$}

\paragraph{Why the leading-order formula is insufficient.}

The accentuation error derived above relies on the approximation
\begin{equation}
  \mathbb{E}_{X,\eta}\!\left[\left(1-R\right)^{2}\right]
  \;\approx\;
  \left(1-R_{\det}\right)^{2},
  \label{eq:acc_leading}
\end{equation}
where $R = \hat{\beta}^{\top}\beta^{*}/\hat{\beta}^{\top}\hat{\beta}$ and
$R_{\det}$ is its deterministic equivalent.  By the second-moment
decomposition,
\begin{equation}
  \mathbb{E}\!\left[(1-R)^{2}\right]
  \;=\;
  \underbrace{(1-\mathbb{E}[R])^{2}}_{\text{bias}^{2}}
  \;+\;
  \underbrace{\operatorname{Var}(R)}_{\text{missing term}}.
  \label{eq:var_decomp}
\end{equation}
The leading-order theory captures $(1-R_{\det})^{2}\approx(1-\mathbb{E}[R])^{2}$
but ignores $\operatorname{Var}(R)$.  When $R\approx 1$
(well-regularized regime), the bias$^{2}$ term is tiny while
$\operatorname{Var}(R)=O(1/n)$ remains finite, so $\operatorname{Var}(R)$
can dominate.

\paragraph{Delta-method setup.}

Define the numerator and denominator of $R$:
\[
  N \;=\; \hat{\beta}^{\top}\beta^{*},
  \qquad
  D \;=\; \hat{\beta}^{\top}\hat{\beta}.
\]
Both $N$ and $D$ concentrate in the large-$n$, fixed-$\gamma=d/n$ limit:
$N\to N_{\det}$, $D\to D_{\det}$, with fluctuations of order $1/\sqrt{n}$.
By the delta method,
\begin{equation}
  R - R_{\det}
  \;\approx\;
  \frac{1}{D_{\det}}\left[(N-N_{\det}) - R_{\det}(D-D_{\det})\right]
  \;=\;
  \frac{\delta Q}{D_{\det}},
  \label{eq:deltaR}
\end{equation}
where $\delta Q = (N - R_{\det}D) - \mathbb{E}[N-R_{\det}D]$.
Because $R_{\det}=N_{\det}/D_{\det}$, the mean of $Q=N-R_{\det}D$ is
zero by construction, so $\delta Q = N - R_{\det}D$ exactly.
Therefore,
\begin{equation}
  \operatorname{Var}(R)
  \;\asymp\;
  \frac{\operatorname{Var}(Q)}{D_{\det}^{2}},
  \quad
  Q \;=\; \hat{\beta}^{\top}(\beta^{*} - R_{\det}\hat{\beta}).
  \label{eq:varR_via_Q}
\end{equation}

\paragraph{Decomposing $\hat{\beta}$ into signal and noise.}

Write $\hat{\beta} = \mathbf{g} + \mathbf{h}$ where
\begin{align*}
  \mathbf{g} &\;=\; (\hat{\Sigma}+\lambda I)^{-1}\hat{\Sigma}\beta^{*}
  \;\asymp\;(\Sigma+\kappa I)^{-1}\Sigma\beta^{*},\\
  \mathbf{h} &\;=\; \sigma(\hat{\Sigma}+\lambda I)^{-1}\!\left(\tfrac{1}{n}X^{\top}\eta\right)
  \;\sim\;
  \mathcal{N}\!\left(0,\,\tfrac{\sigma^{2}}{n}H\right),
  \quad
  H = (\Sigma+\kappa I)^{-1}\Sigma(\Sigma+\kappa I)^{-1}.
\end{align*}
Then
\[
  Q \;=\; (\mathbf{g}+\mathbf{h})^{\top}\bigl(\beta^{*}-R_{\det}(\mathbf{g}+\mathbf{h})\bigr).
\]
Expanding and using $N_{\det}-R_{\det}D_{\det}=0$,
\begin{equation}
  Q
  \;=\;
  \underbrace{\mathbf{h}^{\top}\bigl(\beta^{*}-2R_{\det}\mathbf{g}\bigr)}_{\delta Q^{(\mathrm{lin})}}
  \;-\;
  R_{\det}\!\underbrace{\bigl(\mathbf{h}^{\top}\mathbf{h} - \mathbb{E}[\mathbf{h}^{\top}\mathbf{h}]\bigr)}_{\delta Q^{(\mathrm{quad})}}.
  \label{eq:Q_expand}
\end{equation}
The two terms have zero mean and are uncorrelated (parity argument:
$\delta Q^{(\mathrm{lin})}$ is odd in $\mathbf{h}$, $\delta Q^{(\mathrm{quad})}$
is even), so $\operatorname{Var}(Q)=\operatorname{Var}(\delta Q^{(\mathrm{lin})})+\operatorname{Var}(\delta Q^{(\mathrm{quad})})$.

\paragraph{Linear (noise) contribution.}

Since $\mathbf{h}\sim\mathcal{N}(0,\frac{\sigma^{2}}{n}H)$,
\[
  \operatorname{Var}(\delta Q^{(\mathrm{lin})})
  \;=\;
  \frac{\sigma^{2}}{n}
  \bigl(\beta^{*}-2R_{\det}\mathbf{g}_{\det}\bigr)^{\top}H\,
  \bigl(\beta^{*}-2R_{\det}\mathbf{g}_{\det}\bigr),
\]
with $\mathbf{g}_{\det}=(\Sigma+\kappa I)^{-1}\Sigma\beta^{*}$.
In the eigenbasis of $\Sigma$ (eigenvalues $\lambda_{k}$,
projections $b_{k}=\mathbf{u}_{k}^{\top}\beta^{*}$):
\[
  H_{kk} = \frac{\lambda_{k}}{(\lambda_{k}+\kappa)^{2}},
  \quad
  \bigl(g_{\det}\bigr)_{k} = \frac{\lambda_{k}}{\lambda_{k}+\kappa}b_{k},
  \quad
  w_{k} := \left(\beta^{*}-2R_{\det}\mathbf{g}_{\det}\right)_{k}
  = b_{k}\!\left(1 - \frac{2R_{\det}\lambda_{k}}{\lambda_{k}+\kappa}\right).
\]
Hence,
\begin{equation}
  \operatorname{Var}(\delta Q^{(\mathrm{lin})})
  \;=\;
  \frac{\sigma^{2}}{n}\sum_{k}\frac{\lambda_{k}b_{k}^{2}}{(\lambda_{k}+\kappa)^{2}}
  \!\left(1-\frac{2R_{\det}\lambda_{k}}{\lambda_{k}+\kappa}\right)^{\!2}.
  \label{eq:varQ_lin}
\end{equation}

Expanding the square and introducing shorthand sums:
\begin{align}
  C_{\mathrm{sig}} &:= \sum_{k}\frac{\lambda_{k}}{(\lambda_{k}+\kappa)^{2}}b_{k}^{2},
  \label{eq:Csig}\\
  M_{1} &:= \sum_{k}\frac{\lambda_{k}^{2}}{(\lambda_{k}+\kappa)^{3}}b_{k}^{2},
  \label{eq:M1}\\
  M_{2} &:= \sum_{k}\frac{\lambda_{k}^{3}}{(\lambda_{k}+\kappa)^{4}}b_{k}^{2},
  \label{eq:M2}
\end{align}
equation~\eqref{eq:varQ_lin} becomes
\begin{equation}
  \operatorname{Var}(\delta Q^{(\mathrm{lin})})
  \;=\;
  \frac{\sigma^{2}}{n}\Bigl[C_{\mathrm{sig}} - 4R_{\det}M_{1} + 4R_{\det}^{2}M_{2}\Bigr].
  \label{eq:varQ_lin_compact}
\end{equation}

\paragraph{Quadratic correction.}

$\delta Q^{(\mathrm{quad})} = \mathbf{h}^{\top}\mathbf{h}-\mathbb{E}[\mathbf{h}^{\top}\mathbf{h}]$ is
a centered chi-squared form.  Its variance is
\[
  \operatorname{Var}(\delta Q^{(\mathrm{quad})})
  \;=\;
  2\left(\frac{\sigma^{2}}{n}\right)^{\!2}\operatorname{Tr}[H^{2}]
  \;=\;
  \frac{2\sigma^{4}}{n^{2}}\sum_{k}\frac{\lambda_{k}^{2}}{(\lambda_{k}+\kappa)^{4}}.
  \label{eq:varQ_quad}
\]
This term is $O(\sigma^{4}\gamma/n)$ and dominates when $\sigma^{2}\gg 1/\sqrt{\gamma}$,
i.e.\ at high noise with fixed aspect ratio.

\paragraph{Deterministic equivalent for $\operatorname{Var}(R)$.}

Combining equations~\eqref{eq:varR_via_Q},~\eqref{eq:varQ_lin_compact},
and~\eqref{eq:varQ_quad}:

\begin{equation}
  \boxed{
  \operatorname{Var}(R)
  \;\asymp\;
  \frac{1}{D_{\det}^{2}}
  \left[
    \frac{\sigma^{2}}{n}\Bigl(C_{\mathrm{sig}} - 4R_{\det}M_{1} + 4R_{\det}^{2}M_{2}\Bigr)
    \;+\;
    \frac{2R_{\det}^{2}\sigma^{4}}{n^{2}}\sum_{k}\frac{\lambda_{k}^{2}}{(\lambda_{k}+\kappa)^{4}}
  \right]
  }
  \label{eq:varR_det}
\end{equation}

\paragraph{Full accentuation error including $\operatorname{Var}(R)$.}

The improved deterministic equivalent for the accentuation error is:
\begin{equation}
  \mathbb{E}\!\left[E_{\mathrm{acc}}\right]
  \;\asymp\;
  \beta^{*\top}\Sigma\beta^{*}
  \cdot
  \left[
    (1-R_{\det})^{2}
    \;+\;
    \operatorname{Var}(R)
  \right],
  \label{eq:Eacc_full}
\end{equation}
where $\operatorname{Var}(R)$ is given by~\eqref{eq:varR_det}.

\paragraph{Regime analysis.}
\begin{itemize}
\item \textbf{Low noise} ($\sigma\to 0$): the $\eta$-driven terms
  in~\eqref{eq:varR_det} vanish as $O(\sigma^{2})$, but $\operatorname{Var}(R)$
  remains finite due to \emph{X-fluctuations}
  $\operatorname{Var}_{X}(\mathbb{E}_{\eta}[R|X])=O(1/n)$.
  At low $\sigma$, X-fluctuations dominate
  (empirically $\approx 90\%$ of total $\operatorname{Var}(R)$) and are
  \emph{not} captured by equation~\eqref{eq:varR_det}.
  They require a CLT for quadratic functionals of Wishart matrices
  (a higher-order RMT result).
\item \textbf{Moderate-to-high noise}: the $\eta$-driven terms grow as
  $O(\sigma^{2}/n)$ and quickly dominate ($\approx 97\%$ at $\sigma=1$).
  The quadratic term $O(\sigma^{4}/n)$ can exceed the linear term when
  $\sigma^{2}\sqrt{\gamma}\gg 1$.
\item \textbf{Det.\ equiv.\ accuracy}: even within the $\eta$-driven
  contribution, replacing $H(X)$ by $H_{\det}$ introduces systematic
  error.  When $\kappa=O(\lambda)$ is small relative to the eigenvalue
  spread, the det.\ equiv.\ under-approximates $\operatorname{Tr}[H^{2}]$
  and $\mathbf{w}^{\top}H\mathbf{w}$ (ratio $\sim 0.4$--$0.7$
  in the van Hateren example at $\lambda^{*}\approx 0.001$).
  Using the actual sample eigenvalues of $\hat{\Sigma}$ in place of
  the population ones restores the accuracy (ratio $\to\approx 0.95$).
\end{itemize}

\paragraph{Empirical decomposition (van Hateren $16\times 16$, $d=256$, $n=512$).}

The law of total variance $\operatorname{Var}(R)=\mathbb{E}_{X}[\operatorname{Var}_{\eta}(R|X)]+\operatorname{Var}_{X}(\mathbb{E}_{\eta}[R|X])$ was measured with $T=300$ datasets $\times$ $T'=80$ noise replicas each:

\begin{center}
\begin{tabular}{ccccc}
\hline
$\sigma$ & $\lambda^{*}$ & $\eta$-fraction & X-fraction & formula / $\eta$-MC \\
\hline
0.10 & 0.001 & 10\% & 90\% & 0.69 \\
1.00 & 0.002 & 98\% & 2\%  & 0.44 \\
\hline
\end{tabular}
\end{center}

The formula captures the correct order of magnitude and correctly predicts
that $\operatorname{Var}(R)\sim O(1/n)$, but a complete first-principles
derivation of $\operatorname{Var}(R)$ requires accounting for both
X-fluctuations and the finite-$\kappa$ correction to the det.\ equiv.\
of $H$.
